% 第4章 リターン,リスクと効用(資料4)
\chapter{リターン,リスクと効用}

前章までは将来のキャッシュフローが確定的な世界を扱った.本章からは,将来が
\emph{不確実}な資産(株式など)を扱うための道具立てを導入する.鍵になるのは,
収益率を確率変数とみなす見方,その分布を要約する\textbf{リターン}と
\textbf{リスク},投資家の選好を記述する\textbf{効用関数}と\textbf{期待効用原理},
そしてリスク負担の対価である\textbf{リスク・プレミアム}である.

\section*{導入:株価チャート}

日経平均株価などの株価チャートには,次のような要素が描かれている.

\begin{itemize}
\item \textbf{ローソク足}:一定期間(日足,週足,月足)の\textbf{始値}(はじめね),
\textbf{終値}(おわりね),\textbf{高値}(たかね),\textbf{安値}(やすね)を
1本の「ローソク」で表したもの.
\item \textbf{移動平均線}:過去5日,25日,75日,150日などの終値の平均を
つないだ線.
\item \textbf{出来高(売買高)}:片道数量(売りと買いで一つの取引とした数量)で,
単位は株.ただし,市場が活況か否かという実態は,株数よりも\textbf{売買代金}の方が
よく表現している.大雑把な活況の目安は1日あたり2兆円である.
\end{itemize}

\section{投資収益率}

\subsection{算術収益率}

時刻 $t$ における価格を $P_t$,時刻 $t+1$ における価格を $P_{t+1}$,時刻 $t+1$ に
おける配当を $D_{t+1}$($\geq 0$)とする.

\begin{definition}[算術収益率]
(一期間あたりの)\textbf{算術収益率} $R_{t+1}$ を
\begin{equation}
R_{t+1} = \frac{P_{t+1} - P_t + D_{t+1}}{P_t}
= \frac{P_{t+1} + D_{t+1}}{P_t} - 1
\label{eq:arith-return}
\end{equation}
と定義する.
\end{definition}

分子のうち,$P_{t+1} - P_t$ の部分は,正なら\textbf{キャピタル・ゲイン},負なら
キャピタル・ロスと呼ばれる.$D_{t+1}$ の部分は\textbf{インカム}(インカム・
ゲイン)と呼ばれる.

\subsection{対数収益率}

\begin{definition}[対数収益率]
(一期間あたりの)\textbf{対数収益率} $R^{\log}_{t+1}$ を
\begin{equation}
R^{\log}_{t+1} = \log\frac{P_{t+1}+D_{t+1}}{P_t}
= \log(P_{t+1}+D_{t+1}) - \log P_t
= \log\Bigl(1 + \frac{P_{t+1}-P_t+D_{t+1}}{P_t}\Bigr)
\label{eq:log-return}
\end{equation}
と定義する.
\end{definition}

テーラー展開による近似式 $\log(1+x) \approx x$($|x| \ll 1$ のとき)を用いると,
\begin{equation}
R^{\log}_{t+1} = \log\Bigl(1 + \frac{P_{t+1}-P_t+D_{t+1}}{P_t}\Bigr)
\approx \frac{P_{t+1}-P_t+D_{t+1}}{P_t} = R_{t+1}
\end{equation}
となり,収益率の数値が小さければ対数収益率と算術収益率はほとんど変わらない.

\subsection{対数収益率と算術収益率の違い}

通常,収益率の数値は小さいので両者の値はほとんど変わらないが,\textbf{時系列上の
対称性}に違いがある.

\begin{example}
時系列で価格が 100円 $\to$ 105円 $\to$ 100円と変動したときを考える.
\begin{itemize}
\item 算術収益率では,$5/100 = 5.00\%$,$-5/105 = -4.76\%$ で,平均は $0.12\%$.
\item 対数収益率では,$\log(105/100) = 4.88\%$,$\log(100/105) = -4.88\%$ で,
平均は $0.00\%$.
\end{itemize}
価格が元に戻ったのに算術収益率の平均は正になる一方,対数収益率は(時間に関して)
対称性があり,平均はちょうど0になる.
\end{example}

時系列で収益率を扱う際は,対称性がある方が好ましいと考えられている.ただし,
どちらの方が優れているか,という話ではない.実際,
\begin{itemize}
\item リスク管理では,\textbf{対数収益率}がよく使われる.
\item ポートフォリオの最適化では,\textbf{算術収益率}がよく使われる.
\end{itemize}

\subsection{収益率の分布}

日経平均の日次収益率(1984年1月5日〜2008年11月14日)のヒストグラムを,同じ
平均・分散をもつ正規分布と重ねて描くと,次のことが観察される.

\begin{itemize}
\item 実際の分布は,正規分布に比べて中心が高く尖っている.
\item 裾を拡大して見ると,正規分布ではほとんど起こりえないような大きな下落
(例えば1987年10月20日,ブラックマンデー翌日の約$-15\%$)が実際には起きている.
すなわち,実際の分布は正規分布より\textbf{裾が厚い}(ファット・テイル).
\item 週次収益率(1984年1月4日〜2008年11月10日)でも同様で,2008年10月6日の週
(リーマンショック時)には$-25\%$前後の下落が生じている.
\end{itemize}

このように収益率は日々変動する.そこで,\textbf{収益率は何らかの分布を持つ
確率変数である}と考える.分布を記述する統計量については,
\begin{itemize}
\item 期待値($=$平均値)は水準感の目安になるが,それだけでは情報不足である.
\item 収益率の日々の「変動の大きさ」を示す指標が必要であり,
「分散」または「標準偏差」といった2次の統計量を使用する.
\item 3次以上の統計量(歪度・尖度など)が使われることもある.
\end{itemize}

\section{リターン,リスク}

\subsection{定義}

収益率を確率変数 $R$ とし,観測データを $R_i$, $i=1,2,\dots,N$ とする.

\begin{definition}[リターンとリスク]
\textbf{リターン(期待収益率)}は収益率の期待値 $\E[R]$ であり,観測データからは
\begin{equation}
\E[R] = \frac{1}{N}\sum_{i=1}^{N} R_i
\end{equation}
と推定される.\textbf{リスク}は収益率の標準偏差である.
\end{definition}

なお,「リターン」を(期待値でなく)収益率そのものの意味で使い,期待値の方を
「期待リターン」と言うこともある.人によって定義が違うので注意が必要である.

\subsection{ボラティリティ}

\begin{definition}[ボラティリティ]
収益率の標準偏差を\textbf{ボラティリティ} (volatility) と呼ぶ.省略するときは
「ボラ」または「vol」.
\end{definition}

講義資料では,2012年9月〜2013年6月の日経平均株価について,日々の終値,対数
収益率,および「日々の収益率は独立同一分布に従う」と仮定して過去20日間の収益率
から求めたボラティリティ(年率換算)の時系列が示された.この期間,株価は
約8,800円から15,000円超まで大きく上昇し(いわゆるアベノミクス相場),日次収益率は
$\pm 2\%$程度の変動から,2013年前半には$\pm 6\%$を超える激しい変動へと拡大,
年率換算ボラティリティも15\%程度から45\%超へと大きく変動した.

このように過去データから直接求めたボラティリティを,\textbf{ヒストリカル・
ボラティリティ} (historical volatility) と呼ぶ.
\begin{itemize}
\item データをとる期間は20日間に限らない.
\item データにウェイトをつけること(直近を重視する等)もある.
\end{itemize}
(ヒストリカル・)ボラティリティは重要なリスク尺度の一つであり,
「ボラティリティが高い」$=$「リスクは高い」と考える.

また,利回り(金利)も変動する.1M・6M・12MのLIBORの長期時系列(1996〜2020年頃)
を見ると,金利水準が高い状態と低い状態を行き来しており,こうした動きは
レジーム・スイッチング・モデル (Regime Switching Model) による分析の対象にも
なっている.

実務では,リスク(ボラティリティ)の情報は重視される一方,リターンの情報は
それほど強調されない.リターン(平均)の推定は本質的に難しく,不安定である
ことがその一因である.

\begin{remark}[金融リスク管理におけるリスク]
ボラティリティ,標準偏差,分散などは,\textbf{金融リスク管理では}リスクの大きさを
示す適切な尺度とは考えられていない.なぜならば,これらは平均(あるいは想定水準)
からの乖離の大きさを示すが,リスク管理では想定よりも\emph{儲かる}ことをリスクとは
考えず,\emph{損失のみ}がリスクだからである.そこで,リスク管理向けの(下落量を
示す)尺度が考案されている.例として,VaR (Value at Risk),ES (Expected
Shortfall) などがある.
\end{remark}

\section{効用,効用関数,期待効用原理}

\subsection{どちらを選びますか?}

次の2つの選択問題を考えてみよう.

\begin{enumerate}
\item 安全資産:1年後に\emph{確実に}1万円得られる.危険資産:1年後に20\%の
確率で5万円得られ,80\%の確率で何ももらえない.どちらも5,000円で購入できる.
では,どちらを選ぶ?(期待収益(5,000円)は等しくしてある.)
\item 安全資産:1年後に\emph{確実に}10億円得られる.危険資産:1年後に20\%の
確率で50億円得られ,80\%の確率で何ももらえない.どちらも5億円で購入できる.
では,どちらを選ぶ?(期待収益(5億円)は等しくしてある.)
\end{enumerate}

多くの人は,特に金額が大きい2.では,安全資産を選ぶのではないだろうか.期待収益が
同じでも,不確実な方を避けるこのような選好を,次のように定式化する.

\subsection{リスク回避的投資家}

\begin{keypoint}[リスク回避的投資家]
通常,合理的な投資家は\textbf{リスク回避的}と考える.リスク回避的な投資家は,
\begin{itemize}
\item リターンが同じならば,リスクの小さい方を選ぶ.
\item リスクが同じならば,リターンが大きい方を選ぶ.
\end{itemize}
(ここでは,リターン $=$ 期待収益率.)
\end{keypoint}

なお,別のタイプの投資家として,\textbf{リスク中立的}な投資家,
\textbf{リスク愛好的}な投資家もいる.では,これらの選好を数学的にはどのように
記述するのか.それが以下の効用関数である.

\subsection{効用と効用関数}

$w$ を\textbf{富}または財(金銭・土地・株式・サービス等,価値あるモノの量),
$u$ を\textbf{効用}(富を持つ,またはサービスを受けることで得られる満足度)と
する.

\begin{definition}[効用関数]
次の2つの仮定を必ず満たす関数 $u(w)$ を\textbf{効用関数}という.
\begin{itemize}
\item 仮定1.効用 $u$ は,富 $w$ の関数 $u(w)$ である.
\item 仮定2.富が多いほど,満足度は高い.すなわち
\begin{equation}
u'(w) = \frac{\dd u(w)}{\dd w} > 0.
\end{equation}
\end{itemize}
$u'(w)$ を\textbf{限界効用} (marginal utility) と呼ぶ.
\end{definition}

\subsection{期待効用原理}

\begin{definition}[期待効用原理]
投資家は,自分の効用の数学的期待値 $\E[u(W)]$ を最大化するように行動する
(機会選択を行う).ここで富 $W$ は確率変数である.
\end{definition}

\begin{example}[確率的な富か,期待値相当の確定的な富か]
期待効用原理の下で,投資家は (i) 確率変数の富 $W$,(ii) 期待値相当分の確定的な富
$\E[W]$,のどちらを選択するか?

\textbf{解答.}期待効用原理に従うと,(i) の $\E[u(W)]$ と (ii) の $u(\E[W])$ の
うち,大きい方を選ぶことになる.どちらを選ぶことになるかは $u(w)$ の性質,特に
2階導関数 $\dd^2 u(w)/\dd w^2$ 次第である.
\end{example}

\subsection{Jensenの不等式}

\begin{theorem}[Jensenの不等式]
$f(x)$ が凸関数,$X$ が確率変数のとき,
\begin{equation}
f(\E[X]) \leq \E[f(X)]
\end{equation}
が成り立つ.
\end{theorem}

ここで\textbf{凸関数}とは,$0 < \lambda < 1$ に対して
\begin{equation}
f(\lambda x + (1-\lambda)y) \leq \lambda f(x) + (1-\lambda) f(y)
\end{equation}
が成り立つ関数のことである(内分点の関数値 $\leq$ 関数値の内分点).
\textbf{凹関数}は,$-f(x)$ が凸関数になる関数のことであり,$f$ が凹関数ならば,
Jensenの不等式の不等号の向きは逆になる.

Jensenの不等式は演算の交換に関する不等式の一つであり,ここでは期待値演算と
関数演算の交換に対応している.$X$ が $x_1, x_2$ を各確率0.5でとる場合を図で
考えると,$f(\E[X])$ は曲線上の中点での値,$\E[f(X)]$ は2点の関数値の中点で
あり,凸関数では後者の方が大きいことが図からも見てとれる.

\subsection{効用関数の凹凸とリスク選好}

先の「確率的な富 $W$ か,確定的な富 $\E[W]$ か」の選択は,効用関数の凹凸で
次のように整理できる.

\begin{keypoint}[効用関数の凹凸とリスク選好]
\begin{itemize}
\item $u(w)$ が\textbf{凹関数}($\dd^2u/\dd w^2 < 0$)$\Rightarrow$
$u(\E[W]) \geq \E[u(W)]$ $\Rightarrow$ 確定的な富 $\E[W]$ を選択
(\textbf{リスク回避的}).
\item $u(w)$ が\textbf{凸関数}($\dd^2u/\dd w^2 > 0$)$\Rightarrow$
$u(\E[W]) \leq \E[u(W)]$ $\Rightarrow$ 確率的な富 $W$ を選択
(\textbf{リスク愛好的}).
\item $u(w)$ が\textbf{一次関数}($\dd^2u/\dd w^2 = 0$,凹関数かつ凸関数)
$\Rightarrow$ $u(\E[W]) = \E[u(W)]$ $\Rightarrow$ どちらも選択対象
(\textbf{リスク中立的}).
\end{itemize}
\end{keypoint}

したがって,リスク回避的な投資家の効用関数は
\begin{equation}
u'(w) = \frac{\dd u(w)}{\dd w} > 0 \quad(\text{狭義単調増加関数}),
\qquad
u''(w) = \frac{\dd^2 u(w)}{\dd w^2} < 0 \quad(\text{凹関数})
\end{equation}
という性質をもつ.

\begin{figure}[htbp]
\centering
\begin{tikzpicture}[scale=1.0]
\begin{axis}[
  width=9cm, height=6.5cm,
  xlabel={$w$}, ylabel={$u(w)$},
  xmin=0, xmax=10, ymin=0, ymax=3.6,
  xtick={1,5,9},
  xticklabels={$w_1$,$\E[W]$,$w_2$},
  ytick=\empty,
  axis lines=left]
\addplot[domain=0.3:10, samples=80, thick, blue] {sqrt(x)};
% chord
\addplot[domain=1:9, samples=2, red, dashed] {1 + (3-1)/(9-1)*(x-1)};
% points
\draw[dotted] (axis cs:1,0) -- (axis cs:1,1);
\draw[dotted] (axis cs:9,0) -- (axis cs:9,3);
\draw[dotted] (axis cs:5,0) -- (axis cs:5,2.236);
\draw[dotted] (axis cs:0,2.236) -- (axis cs:5,2.236);
\draw[dotted] (axis cs:0,2) -- (axis cs:5,2);
\node[left] at (axis cs:0,2.236) {\small $u(\E[W])$};
\node[left] at (axis cs:0,2.0) {\small $\E[u(W)]$};
\end{axis}
\end{tikzpicture}
\caption{リスク回避的(凹)な効用関数.富 $W$ が $w_1,w_2$ を各確率0.5でとる
とき,$\E[u(W)]$(赤い弦の中点の高さ)よりも $u(\E[W])$(曲線上の値)の方が
大きい.}
\label{fig:concave-utility}
\end{figure}

\subsection{限界効用逓減の法則}

経済学には\textbf{限界効用逓減(ていげん)の法則}がある.すなわち,富が増える
につれて,1単位の富の追加がもたらす満足度の上昇量は減少していく.これは限界効用
$u'(w)$ が富の単調減少関数であること,すなわち
\begin{equation}
u''(w) = \frac{\dd}{\dd w}\Bigl(\frac{\dd u(w)}{\dd w}\Bigr) < 0
\end{equation}
を意味している.これが成立するのは,リスク回避的な効用関数を用いる場合である.

\subsection{よく使われるリスク回避的な効用関数}

\begin{itemize}
\item \textbf{べき(型)効用関数}:
$u(w) = w^{\gamma}$,\quad $0<\gamma<1$ または $\gamma < 0,\ \gamma\neq 0$.
\item \textbf{指数(型)効用関数}:
$u(w) = -e^{-aw}$,\quad $a>0$.
\item \textbf{対数(型)効用関数}:
$u(w) = \log w$.
\item (※)\textbf{二次効用関数}:
$u(w) = aw - \tfrac{1}{2}bw^2$,\quad $a>0,\ b\geq 0,\ w \leq a/b$.
単調増加する部分のみ使用する.
\end{itemize}

\begin{example}[指数効用による議論の例]
効用関数を $u(w) = -e^{-aw}$($a>0$)とし,富 $W$ の分布は
\[
W = \begin{cases} u & \text{確率 } 0.5,\\ v & \text{確率 } 0.5,\end{cases}
\qquad u < v
\]
とする.投資家は,このまま富 $W$ を保有してもよいし,富の期待値分 $\E[W]$ を
保有することにしてもよい,とする.どちらを選択するか.

$\E[W] = (u+v)/2$ を保有するときの効用は
\[
u(\E[W]) = -e^{-a(u+v)/2},
\]
このまま $W$ を保有したときの期待効用は
\[
\E[u(W)] = -\bigl(e^{-au} + e^{-av}\bigr)/2
\]
である.これらの大小関係をみるために,相加相乗平均の関係式
\[
\frac{a+b}{2} \geq \sqrt{ab},\qquad a>0,\ b>0
\]
を使う.すると,
\[
\E[u(W)] = -\frac{e^{-au}+e^{-av}}{2}
\leq -\sqrt{e^{-au}\times e^{-av}}
= -e^{-a(u+v)/2}
= u(\E[W]).
\]
これより,確実に平均的な富 $\E[W]$ を得られる方を選択する.指数効用では
リスク回避的な投資家の選択と同じになることが確認できた.
\end{example}

\section{リスク・プレミアム}

リスク・プレミアムは,価格にリスクを反映させる際の調整項である.

\subsection{リスク・プレミアムとは}

先の「どちらを選びますか?」の危険資産を,リスク回避的な投資家は購入しない.

\begin{itemize}
\item[問.] では,どうすればリスク回避的な投資家は危険資産でも購入するように
なるのか?
\item[答.] 購入することで抱えるリスクに見合うだけの収益を得られるならば,
危険資産であっても購入する.
\end{itemize}

つまり,危険資産\emph{だから}購入しないのではなく,収益がリスクに見合わないと
考えるから購入しないのである.収益とリスクのバランスが大切である.

\begin{definition}[リスク・プレミアム]
リスクをとること(リスク・テイク)に対する対価(超過収益)を
\textbf{リスク・プレミアム} (risk premium) と呼ぶ.
\end{definition}

では,どれだけの大きさのリスク・プレミアムをとればよいのか.以下では2つの
考え方を示す.

\subsection{確実性等価原理}

定数 $X = 100$,確率変数 $W = 200$(確率50\%)or $0$(確率50\%)とする.
これまでの議論より,リスク回避的投資家では $u(\cdot)$ が凹関数なので
$u(\E[W]) > \E[u(W)]$.ゆえに $X = \E[W]$(安全資産)を選択する.

これを\emph{価格}で考える.安全利子率は $r$(一定)とする.1年後の
キャッシュフロー $X$(確定値)の現在価値は $X/(1+r)$ である.では,1年後の
キャッシュフロー $W$(確率変数)の現在価値はどう評価すべきか.$W$ は確率変数
なので,$X$ とまったく同じ評価をすることはできない.次の2つの考え方がある.

\paragraph{考え方A(割引率の調整項としてのリスク・プレミアム).}
キャッシュフローの期待値 $\E[W]$ を,\textbf{リスク調整済割引率} $R > r$ で
割り引く.すなわち $W$ の現在価値を
\begin{equation}
\text{現在価値} = \frac{\E[W]}{1+R} = \frac{\E[W]}{1 + r + (R-r)}
\end{equation}
とし,差 $R - r$ を\textbf{リスク・プレミアム}と呼ぶ.

これは,学術的な議論よりも実務で時々見られる考え方である.例えば,国が発行する
債券(国債)の利回り $r$ と,一般事業会社が発行する債券(社債)の利回り $R$ の
差 $R-r$ が市場で提示されている.$R-r$(リスク・プレミアム)が高いほどリスクが
高い.例えば,格付が低いほど $R-r$ は高くなる.つまり(同じキャッシュフローに
対して)社債の価格は安くなる.

\paragraph{考え方B(財の調整項としてのリスク・プレミアム).}
効用関数 $u$ を使い,
\begin{equation}
u(w^*) = \E[u(W)]
\label{eq:certainty-equiv}
\end{equation}
を満たす定数 $w^*$ を求め,$W$ の現在価値を $w^*/(1+r)$ とする.つまり,
\begin{enumerate}
\item 効用関数の値が期待効用と同じ値になる富 $w^*$ を求め,
\item その富を\emph{安全利子率} $r$ を使って割り引くことで,現在価値を求める.
\end{enumerate}
この考え方を\textbf{確実性等価原理},$w^*$ を\textbf{確実性等価}と呼ぶ.
このときのリスク・プレミアムは
\begin{equation}
c = \E[W] - w^* = \E[W] - u^{-1}\bigl(\E[u(W)]\bigr)
\end{equation}
であり,富の大きさの低減量(富の大きさへの調整項)である.$c$ はリスク・テイクに
対する対価の大きさを示している(価値を低くすることで,収益率を高めている).

シンプルな確率的な富 $W$($w_1$ または $w_2$ を各確率0.5でとる)で考えると,
凹関数 $u$ のグラフ上で,$\E[u(W)]$ の高さに対応する曲線上の点の横座標が
$w^*$ であり,$w^* < \E[W]$,その差がリスク・プレミアム $c$ である.一般には,
Jensenの不等式より $u(\E[W]) \geq \E[u(W)]$ かつ $u(\cdot)$ は単調増加なので,
$u(\E[W]-c) = \E[u(W)]$ を満たす $c$ は $c \geq 0$ となる.

\begin{example}[例題4-1]
現在の富を100万円とし,価格 $p$ 万円の宝くじの購入を検討する.その宝くじは,
確率50\%で100万円,確率50\%で何ももらえない.安全利子率を $r$(年率),宝くじの
支払いまでの期間を $\Delta t$ 年とする.効用関数は $u(w) = w^{0.5}$($w$:万円)
とする.

宝くじを「購入する」「購入しない」の選択肢がある.
\begin{itemize}
\item 購入しないときの効用は(現在の富を安全に運用すると考えて)
\[
u\bigl(100\times(1+r\Delta t)\bigr) = \bigl(100(1+r\Delta t)\bigr)^{0.5}.
\]
\item 購入するときの期待効用は(宝くじ購入の支出,残金の安全な運用,宝くじの
結果から)
\[
0.5\bigl((100-p)(1+r\Delta t) + 100\bigr)^{0.5}
+ 0.5\bigl((100-p)(1+r\Delta t) + 0\bigr)^{0.5}.
\]
\end{itemize}
以上より,宝くじを購入すべき条件は
\[
0.5\bigl((100-p)(1+r\Delta t)+100\bigr)^{0.5}
+ 0.5\bigl((100-p)(1+r\Delta t)\bigr)^{0.5}
\geq \bigl(100(1+r\Delta t)\bigr)^{0.5}.
\]

簡単のため $\Delta t = 0$(即座に結果がわかる)とすると,この条件は
\[
0.5\sqrt{200-p} + 0.5\sqrt{100-p} \geq 10
\qquad\Longleftrightarrow\qquad
p \leq 43.75 \text{ (万円)}
\]
となる.

次に,$p = 43.75$ 万円であったとする.このとき,購入しないときの富の効用と,
購入時の期待効用は等しい.そのため,購入時の期待効用の値を与える確定的な富
$w^*$(確実性等価)がわかる:$w^* = 100$(購入しないときの富だから)である.
一方,購入時の富の期待値は($\Delta t = 0$ を使用)
\[
\E[W] = \bigl(\text{初期富}(100) - \text{宝くじ価格}(43.75)\bigr) +
\text{宝くじの期待値}(50) = 56.25 + 50 = 106.25.
\]
よって,財の調整項であるリスク・プレミアム $c$ は
\[
c = \E[W] - w^* = 106.25 - 100 = 6.25 \text{ (万円)}.
\]
\end{example}

\begin{example}[例題4-2]
例題4-1と同じ設定で,現在の富だけを50万円に変える.購入しないときの効用は
$(50(1+r\Delta t))^{0.5}$,購入するときの期待効用は
\[
0.5\bigl((50-p)(1+r\Delta t)+100\bigr)^{0.5}
+ 0.5\bigl((50-p)(1+r\Delta t)\bigr)^{0.5}
\]
であり,購入すべき条件は($\Delta t = 0$ として)
\[
0.5\sqrt{150-p} + 0.5\sqrt{50-p} \geq \sqrt{50}
\qquad\Longleftrightarrow\qquad
p \leq 37.5 \text{ (万円)}.
\]
$p = 37.5$ 万円のとき,確実性等価は $w^* = 50$(購入しないときの富),
$\E[W] = (50 - 37.5) + 50 = 62.5$ なので,リスク・プレミアムは
\[
c = \E[W] - w^* = 62.5 - 50 = 12.5 \text{ (万円)}.
\]
\end{example}

例題4-1と4-2を比べると,\textbf{購入判断の分岐となる宝くじ価格 $p$ も
リスク・プレミアム $c$ も,現在の富に依存}することがわかる(富が少ないほど,
同じ宝くじに払える価格は低く,要求するリスク・プレミアムは大きい).では,
\emph{全員にとって}合理的な宝くじの価格 $p$,リスク・プレミアム $c$ はあるの
だろうか?

\subsection{効用関数,分布の変換,価格}

例題4-1での投資 $W$ の価値は $w^*/(1+r\Delta t) = (\E[W]-c)/(1+r\Delta t) = 100$,
例題4-2では同様に $50$ である.価格を知るには $w^*$ が必要である.例題では
$w^*$ を所与としたが,一般に $w^*$ は $u(w)$ を使って数値計算で求めることになり,
これは大変である.

ここで,もしも $u(w) = w$(リスク中立)ならば,
$w^* = u^{-1}(\E[u(W)]) = \E[W]$ となるから,$w^*$ は簡単に計算できることに
注目する.そこで,次のような発想の転換を行う.

\begin{keypoint}[リスク中立化の発想(preview)]
投資家はリスク中立的($u(w)=w$)とみなして評価するが,それでは現実と異なるので,
辻褄のあうように $W$ の\emph{分布}を変更する.つまり,
\[
(\text{実際の } u(w),\ \text{実際の } W \text{ の分布})
\]
を,それと同じ価格を与える
\[
(\text{リスク中立な } u(w)=w,\ \text{リスク中立な } W \text{ の分布})
\]
に変換して計算する.そのためには,実際のデータからリスク中立な $W$ の分布を
求める必要がある.
\end{keypoint}

これはかなり advanced な話であるが,後の章で学ぶデリバティブの価格付け
(リスク中立化法)の基本発想がここに現れている.

\subsection{Sharpe Ratio(シャープレシオ)}

ここから先は実務的な話である.「リスクに見合うリスク・プレミアムがついているか」
が金融商品の評価の尺度になる.考え方A(割引率の調整項としてのリスク・プレミアム)
で考える.

\begin{definition}[シャープレシオ]
ある資産の期待収益率を $\mu$,ボラティリティを $\sigma$,安全利子率を $r$ と
する.このとき
\begin{equation}
SR = \frac{\mu - r}{\sigma}
\end{equation}
を \textbf{Sharpe Ratio}(シャープ・レシオ,SR)という.
\end{definition}

これは「超過収益率 $\div$ ボラティリティ(orリスク)」であり,その意味は
\textbf{リスク一単位当たりに要求される超過収益率}である.実務では,投資商品の
パフォーマンスを示す尺度として使われる.

\subsection{ファンドのパフォーマンス評価}

\begin{definition}[ファンド]
\textbf{ファンド}とは,投資のために集めた資金,またはその資金による投資運用商品
のこと.
\begin{itemize}
\item \textbf{公募ファンド}:不特定多数からオープンに募集するファンド.投資信託等.
\item \textbf{私募ファンド}:限られた人から集めるファンド.ヘッジファンド等.
\end{itemize}
\end{definition}

SRによるファンドのパフォーマンス評価は次のように行う.
\begin{itemize}
\item 過去の価格データを使用する.
\item 一定期間の収益率,収益率の標準偏差(リスク)を計算し,SRを算出する.
\item SRが高いほどパフォーマンスが良い,と考える(ただし $SR>0$ の場合).
\end{itemize}
なお,SRが良くなくても敢えて投資することもある(分散投資の効果.後述の章で
扱う).

\begin{example}[例題4-3]
資産Aの収益率は,確率50\%で25\%,確率50\%で$-15\%$である.資産Bの収益率は,
確率50\%で30\%,確率50\%で$-20\%$である.安全利子率を2\%として,それぞれの
資産のシャープレシオを求め,どちらの方が投資効率が良いか議論しなさい.

\medskip
\noindent\textbf{解答.}
資産Aは,リターン $0.5\times 25 + 0.5\times(-15) = 5\%$,リスク(収益率の標準
偏差)は平均5\%からの乖離が $\pm 20\%$ なので20\%である.よって
\[
SR_A = \frac{5-2}{20} = 0.15.
\]
資産Bは,リターン $5\%$,リスク25\%なので,
\[
SR_B = \frac{5-2}{25} = 0.12.
\]
SRは資産Aの方が高いので,SRの観点からは資産Aの方が投資効率が良い
(パフォーマンスが良い)と考えられる.
\end{example}

\subsection{インデックス,トラッキング・エラー,IR}

ファンドの収益率だけで評価してよいのだろうか.頑張って成果を出しているファンド
でも,市場全体が落ち込んでいく局面ではどうにもならない.一方で,たいした成果を
出していないように見えるファンドでも,市場全体が上がっているときは Sharpe Ratio
は高くなる.そこで,市場全体の動きを補正して考えることが行われる.

\begin{itemize}
\item \textbf{インデックス}:市場平均価格.市場価格の動向を示す指数.
例:日経225,TOPIX,NYダウ(米),FT100(フッツィー100,英)など.
\item \textbf{パッシブ運用(インデックス運用)}:インデックスの動きに沿った運用.
対義語は\textbf{アクティブ運用}.ちなみに,アクティブファンドは手数料が高めである.
\item \textbf{アクティブ・リターン} $\mu'$:ファンドとインデックスの収益率の差.
収益率そのものでなく,アクティブ・リターンで評価する.
\item $\E[\mu']$:アクティブ・リターンの期待値.
\item $\sigma'$:アクティブ・リターンの標準偏差.これを
\textbf{トラッキング・エラー}と呼ぶ.
\end{itemize}

\begin{definition}[インフォメーション・レシオ]
\begin{equation}
IR = \frac{\E[\mu']}{\sigma'}
\end{equation}
を\textbf{インフォメーション・レシオ} (Information Ratio, IR) という.
\end{definition}

現在は,$\sigma'$(トラッキング・エラー)や IR の方が評価指標としてよく
使われている.

% ---------------- 演習問題 ----------------
\section{演習問題}

\begin{exercise}[問題4-1]
現在の富を100万円とし,価格 $p$ 万円の宝くじの購入を検討する.その宝くじは,
確率5\%で500万円,確率95\%で10万円もらえるものとする.効用関数は
$u(w) = w^{0.5}$($w$:万円)とする.
\begin{enumerate}
\item[(1)] 購入が合理的と考えられるときの宝くじの価格 $p$ の条件を求めなさい.
\item[(2)] 価格 $p$ に対する確実性等価が100万円(投資をしないときと同じ)で
あったとする.このときの宝くじのリスク・プレミアムを求めなさい.
\end{enumerate}
解答は万円単位で,小数第2位まで求めること.なお,宝くじの購入も当選額の支払いも
即座に行われるとする(例題4-1で $\Delta t = 0$ の意味).
\end{exercise}

\begin{answerbox}
\textbf{考え方.}例題4-1と同じ枠組みである.「購入しない場合の効用」と「購入する
場合の期待効用」を比較し,後者が前者以上になる価格の範囲を求める.リスク・
プレミアムは $c = \E[W] - w^*$ で計算する.

\medskip
\textbf{(1)}
宝くじを購入しないときの効用は $u(100\times(1+r\Delta t))$.購入するときの
期待効用は
\[
0.05\bigl((100-p)(1+r\Delta t)+500\bigr)^{0.5}
+ 0.95\bigl((100-p)(1+r\Delta t)+10\bigr)^{0.5}.
\]
購入が合理的となるための条件は
\[
0.05\bigl((100-p)(1+r\Delta t)+500\bigr)^{0.5}
+ 0.95\bigl((100-p)(1+r\Delta t)+10\bigr)^{0.5}
\geq u\bigl(100(1+r\Delta t)\bigr).
\]
$\Delta t = 0$ とすると,この条件は
\[
0.05\sqrt{600-p} + 0.95\sqrt{110-p} \geq \sqrt{100} = 10
\]
となり,数値的に解くと
\[
\boxed{p \leq 24.19 \text{ (万円)}}
\]
を得る.$p = 24.19$ 万円のとき,購入する場合と購入しない場合の期待効用は等しく,
どちらも選択しうる状態になっている.

\medskip
\textbf{(2)}
$p = 24.19$ 万円のときは,宝くじを購入しない場合の富($=$初期富100万円)が
確実性等価となる($w^* = 100$).確実性等価における効用は,宝くじを購入する
場合の期待効用に等しい(購入するもしないも,期待効用の面からは同様に合理的で
ある状態).

宝くじの期待値は $0.05\times 500 + 0.95\times 10 = 25 + 9.5 = 34.5$ 万円なので,
\[
\E[W] = \text{初期富}(100) - \text{宝くじ価格}(24.19) + \text{宝くじの期待値}(34.5)
= 110.31 \text{ 万円}.
\]
よって,リスク・プレミアムは
\[
c = \E[W] - w^* = 110.31 - 100 = \boxed{10.31 \text{ (万円)}}.
\]

\medskip
\textbf{ポイント.}この宝くじの「期待値」は34.5万円だが,リスク回避的な投資家
(平方根型効用)が支払ってもよい上限は24.19万円にとどまる.その差がリスクを
嫌うことによる値引き分であり,確実性等価の言葉で言えばリスク・プレミアム
$c=10.31$ 万円に対応する.
\end{answerbox}

\begin{exercise}[問題4-2]
例題4-3に,さらに資産Cが選択肢として加わった.資産Cの収益率は,確率60\%で
15\%,確率40\%で$-12\%$である.安全利子率を2\%とする.
\begin{enumerate}
\item[(1)] 資産Cのシャープレシオを求めなさい.
\item[(2)] どの資産への投資が効率が良いか,それはなぜかをSR(シャープレシオ)の
数値から議論しなさい.
\end{enumerate}
\end{exercise}

\begin{answerbox}
\textbf{考え方.}2値をとる収益率の期待値と標準偏差を求め,$SR=(\mu-r)/\sigma$ に
代入する.

\medskip
\textbf{(1)}
資産Cのリターンは
\[
\mu_C = 0.6\times 15 + 0.4\times(-12) = 9 - 4.8 = 4.2\%.
\]
分散は
\[
\sigma_C^2 = 0.6\times(15-4.2)^2 + 0.4\times(-12-4.2)^2
= 0.6\times 10.8^2 + 0.4\times 16.2^2 = 174.96
\]
より,リスクは $\sigma_C = 13.227\ldots\%$.よって
\[
SR_C = \frac{4.2 - 2}{13.227\ldots} = 0.1663\ldots.
\]

\medskip
\textbf{(2)}
3資産のSRを並べると,
\[
SR_A = 0.15,\qquad SR_B = 0.12,\qquad SR_C = 0.166.
\]
SRは資産Cが最も高いので,資産Cの投資効率が最も高い(パフォーマンスが良い)と
考えられる.資産Cは,リターン(4.2\%)は資産AとB(いずれも5\%)より低いが,
リスク(13.2\%)が大幅に低いため,リスク一単位当たりの超過収益率であるSRで
測るとパフォーマンスは高いと評価される.

\medskip
\textbf{注意.}パフォーマンス評価の尺度はSRだけではない.また,SRによる評価が
必ずしも適切とは限らない(IR (Information Ratio) などもある).
\end{answerbox}
